
\documentclass{article}
\usepackage{colortbl}
\usepackage{makecell}
\usepackage{multirow}
\usepackage{supertabular}

\begin{document}

\newcounter{utterance}

\centering \large Interaction Transcript for game `referencegame', experiment `letter\_grids', episode 8 with reflexion{-}Qwen3.5{-}4B.
\vspace{24pt}

{ \footnotesize  \setcounter{utterance}{1}
\setlength{\tabcolsep}{0pt}
\begin{supertabular}{c@{$\;$}|p{.15\linewidth}@{}p{.15\linewidth}p{.15\linewidth}p{.15\linewidth}p{.15\linewidth}p{.15\linewidth}}
    \# & $\;$A & \multicolumn{4}{c}{Game Master} & $\;\:$B\\
    \hline

    \theutterance \stepcounter{utterance}  
    & & \multicolumn{4}{p{0.6\linewidth}}{
        \cellcolor[rgb]{0.9,0.9,0.9}{
            \makecell[{{p{\linewidth}}}]{
                \texttt{\tiny{[P1$\langle$GM]}}
                \texttt{You have attempted this same task before and did not succeed. These are the plans you wrote after those attempts.} \\
\texttt{Trial \#1: New plan: My previous expression failed because it was too vague and did not explicitly distinguish the target grid from the distractors. Specifically, I described the top and bottom rows and the first column, which are common to all three grids, but I failed to mention the unique pattern in the middle rows: the target has only the first column filled in the middle, whereas Distractor 1 has alternating columns and Distractor 2 has three filled columns. Next time, I will construct the expression by first stating the common features (solid top/bottom, left column) and then immediately contrasting the middle rows by specifying that the target has "only the first column filled" while the others have different patterns, ensuring the description uniquely identifies the target.} \\
\\ 
\texttt{Trial \#2: New plan: My previous expression failed because it was too generic and did not explicitly contrast the middle rows of the target grid against the specific patterns of the distractors. Specifically, I mentioned the solid top/bottom rows and the left column, which are common to all grids, but I failed to clearly distinguish the middle rows: the target has only the first column filled in the middle, whereas Distractor 1 has alternating filled cells and Distractor 2 has three filled columns. Next time, I will construct the expression by first stating the common features (solid top/bottom rows and a solid left column) and then immediately contrasting the middle rows by specifying that the target has "only the first column filled" while the others have "alternating" or "three filled columns" patterns, ensuring the description uniquely identifies the target.} \\
\\ 
\\ 
\texttt{You are given three grids, where each of them is 5 by 5 in size.} \\
\texttt{Grids have empty cells marked with "▢" and filled cells marked with "X".} \\
\texttt{Your task is to generate a referring expression that best describes the target grid while distinguishing it from the two other distractor grids.} \\
\texttt{The first grid is the target grid, and the following two grids are the distractors.} \\
\\ 
\texttt{Target grid:} \\
\\ 
\texttt{X X X X X} \\
\texttt{X □ □ □ □} \\
\texttt{X X X X X} \\
\texttt{X □ □ □ □} \\
\texttt{X X X X X} \\
\\ 
\texttt{Distractor grid 1:} \\
\\ 
\texttt{X X X X X} \\
\texttt{X □ X □ X} \\
\texttt{X X X X X} \\
\texttt{X □ X □ X} \\
\texttt{X X X X X} \\
\\ 
\texttt{Distractor grid 2:} \\
\\ 
\texttt{X X X X X} \\
\texttt{X □ □ □ □} \\
\texttt{X □ □ □ □} \\
\texttt{X □ □ □ □} \\
\texttt{X X X X X} \\
\\ 
\texttt{Instruction: Describe the target grid.} \\
\texttt{Generate the referring expression starting with the tag "Expression: " for the given target grid. Omit any other text.} \\
            }
        }
    }
    & \\ \\

    \theutterance \stepcounter{utterance}  
    & \multicolumn{4}{p{0.6\linewidth}}{
        \cellcolor[rgb]{0.9,0.9,0.9}{
            \makecell[{{p{\linewidth}}}]{
                \texttt{\tiny{[P1$\rangle$GM]}}
                \texttt{Expression: The grid with a solid top row, a solid bottom row, and a solid left column, where the middle rows contain only the first column filled.} \\
            }
        }
    }
    & & \\ \\

    \theutterance \stepcounter{utterance}  
    & & & \multicolumn{2}{p{0.3\linewidth}}{
        \cellcolor[rgb]{0.9,0.9,0.9}{
            \makecell[{{p{\linewidth}}}]{
                \texttt{\tiny{[GM$|$GM]}}
                \texttt{Expression: The grid with a solid top row, a solid bottom row, and a solid left column, where the middle rows contain only the first column filled.} \\
            }
        }
    }
    & & \\ \\

    \theutterance \stepcounter{utterance}  
    & & \multicolumn{4}{p{0.6\linewidth}}{
        \cellcolor[rgb]{0.9,0.9,0.9}{
            \makecell[{{p{\linewidth}}}]{
                \texttt{\tiny{[GM$\rangle$P2]}}
                \texttt{You have attempted this same task before and did not succeed. These are the plans you wrote after those attempts.} \\
\texttt{Trial \#1: New plan: I incorrectly identified the target grid by misinterpreting the description of the middle rows. The expression specified that only the **first column** is filled in the middle rows, but I selected the grid where the **first two columns** are filled. Next time, I will carefully verify that the filled cells in the middle rows align strictly with the first column (column index 1) and exclude any filled cells in the second column. I will also double{-}check that the top and bottom rows are fully filled before making my final selection.} \\
\\ 
\texttt{Trial \#2: New plan: I misinterpreted the description of the "middle row" in the expression. The prompt specified a "solid middle row," which implies the entire third row is filled with "X". However, I likely focused only on the "two rows in between" containing only the first column filled, causing me to overlook the requirement for the third row to be completely solid. Next time, I will strictly verify that the third row contains five "X"s (a solid row) in addition to the top, bottom, and middle{-}left columns being filled, ensuring no gaps exist in the third row before selecting the grid.} \\
\\ 
\\ 
\texttt{You are given three grids, where each of them is 5 by 5 in size.} \\
\texttt{Grids have empty cells marked with "▢" and filled cells marked with "X".} \\
\texttt{You are also given a referring expression that describes one of the given grids.} \\
\texttt{Your task is to select the grid that matches the given referring expression.} \\
\texttt{Generate only the number (in text) of the grid that matches the given expression by selecting first, second, or third.} \\
\\ 
\texttt{First:} \\
\\ 
\texttt{X X X X X} \\
\texttt{X □ □ □ □} \\
\texttt{X □ □ □ □} \\
\texttt{X □ □ □ □} \\
\texttt{X X X X X} \\
\\ 
\texttt{Second:} \\
\\ 
\texttt{X X X X X} \\
\texttt{X □ X □ X} \\
\texttt{X X X X X} \\
\texttt{X □ X □ X} \\
\texttt{X X X X X} \\
\\ 
\texttt{Third:} \\
\\ 
\texttt{X X X X X} \\
\texttt{X □ □ □ □} \\
\texttt{X X X X X} \\
\texttt{X □ □ □ □} \\
\texttt{X X X X X} \\
\\ 
\texttt{Expression: The grid with a solid top row, a solid bottom row, and a solid left column, where the middle rows contain only the first column filled.} \\
\texttt{Question: Which grid does the expression refer to?} \\
\texttt{Start with the tag "Answer: ", followed by your selection. Omit any other text.} \\
            }
        }
    }
    & \\ \\

    \theutterance \stepcounter{utterance}  
    & & & \multicolumn{4}{p{0.6\linewidth}}{
        \cellcolor[rgb]{0.9,0.9,0.9}{
            \makecell[{{p{\linewidth}}}]{
                \texttt{\tiny{[GM$\langle$P2]}}
                \texttt{Answer: First} \\
            }
        }
    }
     \\ \\

    \theutterance \stepcounter{utterance}  
    & & & \multicolumn{2}{p{0.3\linewidth}}{
        \cellcolor[rgb]{0.9,0.9,0.9}{
            \makecell[{{p{\linewidth}}}]{
                \texttt{\tiny{[GM$|$GM]}}
                \texttt{first} \\
            }
        }
    }
    & & \\ \\

    \theutterance \stepcounter{utterance}  
    & \multicolumn{4}{p{0.6\linewidth}}{
        \cellcolor[rgb]{0.95,0.95,0.95}{
            \makecell[{{p{\linewidth}}}]{
                \texttt{\tiny{[P2$\rightarrow$P1]}}
                \texttt{Post{-}episode Reflexion (not sent to GM):} \\
\texttt{New plan: My previous attempts failed because the descriptions of the middle rows were too vague ("different patterns" or "alternating") and did not explicitly state the exact count of filled cells in the middle rows for each grid. The target grid has exactly one filled cell in the middle rows (the first column), while Distractor 1 has two and Distractor 2 has three. Next time, I will construct the expression by first identifying the common solid top/bottom rows and left column, then immediately specifying that the middle rows contain "exactly one filled cell in the first column" to uniquely distinguish it from the distractors which have two or three filled cells respectively.} \\
            }
        }
    }
    & & \\ \\

    \theutterance \stepcounter{utterance}  
    & & & \multicolumn{4}{p{0.6\linewidth}}{
        \cellcolor[rgb]{0.95,0.95,0.95}{
            \makecell[{{p{\linewidth}}}]{
                \texttt{\tiny{[P1$\rightarrow$P2]}}
                \texttt{Post{-}episode Reflexion (not sent to GM):} \\
\texttt{New plan: I misinterpreted the phrase "middle rows" as referring to the two rows between the top and bottom, but the expression actually describes a specific pattern where the **third row** (the true middle) must be a solid line of "X"s, while the rows immediately above and below it contain only the first column filled. The target grid (Second) has a solid third row, whereas the First and Third grids have gaps in the third row. Next time, I will strictly check if the third row is completely filled with "X" before considering the side columns, ensuring I do not confuse the "middle rows" plural with the single solid middle row.} \\
            }
        }
    }
     \\ \\

\end{supertabular}
}

\end{document}
